% !PRINTED 621
\weberpage{621}
%<<CONT
and herein the coefficients
\[
f(1), f'(1), \tfrac12 f''(1), \cdots, \dfrac{1}{\Pi(\mu-1)} f^{(\mu-1)}(1)
\]
are integral rational numbers, namely
\begin{align*}
\dfrac{1}{\Pi(\mu-1)} f^{(\mu-1)}(1) &= x_{\mu-1} \\[6pt]
\dfrac{1}{\Pi(\mu-2)} f^{(\mu-2)}(1) &= x_{\mu-2}+(\mu-1)x_{\mu-1} \\[6pt]
\dfrac{1}{\Pi(\mu-3)} f^{(\mu-3)}(1) &= x_{\mu-3}+(\mu-2)x_{\mu-2}+\dfrac{(\mu-1)(\mu-2)}{1.2}x_{\mu-1} \tag{5}\\[6pt]
\cdots\cdots\cdots &\cdots\cdots\cdots\cdots\cdots\cdots\cdots \\[6pt]
f(1) &= x_0+x_1+\cdots\cdots\cdots\cdots+x_{\mu-1}.
\end{align*}
Now since $\omega$ is divisible by $q$, and consequently by $\sigma^\mu$,
$f(1)$ must by (4) be divisible by $q$. But since $\mu>1$, and
consequently $q:\sigma$ is still divisible by $\sigma$, it follows that
\[
f'(1) - \dfrac{\sigma}{1.2} f''(1) \ldots
\]
also is divisible by $\sigma$, and consequently $f'(1)$ by $q$; and reasoning
onward in this way, it results that all the rational
numbers (5) are divisible by $q$, and consequently the
$x_{\mu-1}, x_{\mu-2}, \ldots, x_0$ are also divisible by $q$. With this Theorem II. is completely proved.

Instead of the sequence (1) one may also take as a basis of $\mathfrak{o}$ an
arbitrary sequence of $\mu$ successive powers of $r$:
\weq{6}{r^a, r^{a+1}, r^{a+2}, \ldots, r^{a+\mu-1},}
which results when one multiplies all the terms of the sequence (1)
by $r^a$.

The correctness of this remark follows from the fact that
the $r^{\pm a}$ are integers, and that therefore $\omega$ and $r^a\omega$ are always simultaneously integral or fractional numbers.

For example, the numbers
\weq{7}{r^{-\frac12\mu+1}, r^{-\frac12\mu+2}, \ldots, r^{-1}, 1, r, r^2, \ldots, r^{\frac12\mu}}
form a basis of $\mathfrak{o}$.

With this the fundamental number $\Delta$ of the field $\Omega_m$ is also determined.
By the general definition (\S.\,145) it is equal to the discriminant % [sic] German source prints "Discirminante"; "Discriminante" is meant
\[
\Delta(1, r, r^2, \ldots, r^{\mu-1}),
\]
i.e. equal to the discriminant of the cyclotomic equation determined in \S.\,169.
